% 19 November 2026. Drafted for Edward Baker with OpenAI Codex assistance.
% Lead actual backend/effort unavailable; helpers explicitly Sol6.1/ultra.
\section{Finite-cell propagation and a sharper complete allowance}
\label{sec:local-estimates}

The first continuation fixes the physical reference $P_N$ of
\eqref{as:PN} and compares two ways to acquire its residual on one
counted contour. The signed endpoint identity \eqref{prog:centered-error}
is inherited. The new estimates below fill in the lateral boundary
payment and sharpen the complete value-to-discriminant allowance. They
are finite estimates, not an RH-equivalent reformulation. Their arithmetic
application must still acquire the count, clearance and positive margin.

\subsection{One fixed arithmetic cell}

Let $0<U<t_*\le1/2$, $0<h\le1$, $X-A\ge200$, and require
\begin{equation}
 N^2\le\frac{X-A}{4\pi}+\frac{t_*-U}{16},\qquad
 \frac{X+A}{4\pi}+\frac{t_*}{16}<(N+1)^2.
 \label{le:cutoff}
\end{equation}
The natural cutoff is then the same integer throughout
$0\le u\le U$, $|\Re z-X|\le A$, $|\Im z|\le h$, with
$t=t_*-u$. This avoids a seam payment without discarding arithmetic terms.
It permits an interval of centers precisely when
$2A<4\pi(2N+1)-\pi U/4$, before imposing $X-A\ge200$.
Use the fixed circle $\Gamma:|z-X|=r<\min(A,h)$.
Neither its two-zero count nor a real positive reference discriminant is
assumed to follow from the cutoff inequalities.

For clarity, the complete published value budget can be written without
the normalizer. Put $T=x/2$, $\eta=|y|$, $\sigma_\pm=(1\pm\eta)/2$,
$w_\pm=\sigma_\pm+\ii T$, $T'=T+\pi t/8$, and
$a_{\rm RS}=\sqrt{T'/(2\pi)}$. Define
\begin{align}
 \epsilon_n(t,w)&=\exp\left(
 \frac{(t^2/8)|\alpha(w)-\log n|^2+t/4+1/6}{T-3.33}
 \right)-1,\notag\\
 \delta_\sigma(t,T)&=
 \left(\frac{0.397\,9^\sigma}{a_{\rm RS}-0.865}
                 +\frac5{3(T-6)}\right)
                 \exp\left(\frac{3.49}{T-4}\right).
 \label{le:source-errors}
\end{align}
Then Corollary~6.5 of \cite{Polymath2019} gives
\begin{align}
 |H_t(z)-P_N(t,z)|\le\mathcal V_N(t,x,\eta)
 :={}&\sum_{n=1}^N\sum_{\nu\in\{-,+\}}
                 |p_n(t,w_\nu)|\epsilon_n(t,w_\nu)\notag\\
 &+e^{t\pi^2/64}|\mathfrak M_0(\ii T')|
                   (1+\delta_{\sigma_-}+\delta_{\sigma_+}).
 \label{le:source-value}
\end{align}
The $1$ in the second line pays the leading omitted correction. It is
not replaced by the smaller retained-correction error. The source domain
is $0<t\le1/2$, $x\ge200$, $|y|\le1$ and $N=\lfloor a_{\rm RS}\rfloor$;
\eqref{le:cutoff} enforces this throughout the cell. The source's branch
names differ from ours; the explicit upper-half-plane arguments $w_\pm$
and conjugation fix that convention. The formula supplies initial and
lateral data as well as a direct endpoint envelope. The full finite
forcing remains \eqref{rt:physical-defect}, with its complex sum retained
before taking a modulus whenever it is evaluated jointly.

\subsection{Bounded-domain heat propagation with both boundaries}

Write $x_-=X-A$, $x_+=X+A$, $I=(x_-,x_+)$ and $L=2A$.
For every fixed $y\in[-h,h]$, $e(u,x)=E(u,x+\ii y)$ satisfies
$e_u=e_{xx}+\mathcal R_N$. Suppose the data and forcing are continuous,
and $E$ is $C^1$ in time, $C^2$ in $x$ and holomorphic on a neighborhood
of the closed rectangle, locally uniformly in time. These hypotheses
hold for the fixed finite reference and canonical field on this cell.
There is no whole-line extension premise.

Let $K(\tau;x,\xi)$ be the Dirichlet heat kernel on $I$:
\begin{equation}
 K(\tau;x,\xi)=\frac2L\sum_{n\ge1}e^{-(n\pi/L)^2\tau}
  \sin\frac{n\pi(x-x_-)}L\sin\frac{n\pi(\xi-x_-)}L,
 \qquad \tau>0.
 \label{le:kernel}
\end{equation}
It is nonnegative. Put
$k_-=\partial_\xi K|_{\xi=x_-}$ and
$k_+=-\partial_\xi K|_{\xi=x_+}$, both nonnegative. With
$f_y(\xi)=E(0,\xi+\ii y)$ and $g_\pm(v,y)=E(v,x_\pm+\ii y)$,
the exact formula is
\begin{align}
 E(u,x+\ii y)={}&\int_I K(u;x,\xi)f_y(\xi)\dd\xi\notag\\
 &+\int_0^u\int_I K(u-v;x,\xi)
             \mathcal R_N(t_*-v,\xi+\ii y)\dd\xi\dd v\notag\\
 &+\int_0^u\{k_-(u-v;x)g_-(v,y)
                    +k_+(u-v;x)g_+(v,y)\}\dd v.
 \label{le:bounded-duhamel}
\end{align}
Indeed, differentiate $\int_I K(u-v;x,\xi)E(v,\xi+\ii y)\dd\xi$.
Green's identity yields the boundary term
$[K E_\xi-K_\xi E]_{x_-}^{x_+}$ and the forcing integral. Integration
in $v$, followed by the approximate-identity limit at $u$, proves
\eqref{le:bounded-duhamel}. The variable $y$ is a parameter, so no top
or bottom boundary integral enters this one-dimensional heat problem.

Set $S(\tau,x)=\int_I K(\tau;x,\xi)\dd\xi$ and
$Q_\pm(\tau,x)=\int_0^\tau k_\pm(s;x)\dd s$. Applying the formula to
the constant solution gives
\begin{equation}
 S+Q_-+Q_+=1.
 \label{le:mass}
\end{equation}
The maximum principle, comparing the boundary responses with the
half-line heat solution, gives
\begin{equation}
 1-S(\tau,x)\le\min\left\{1,
 \operatorname{erfc}\frac{x-x_-}{2\sqrt\tau}
 +\operatorname{erfc}\frac{x_+-x}{2\sqrt\tau}\right\}.
 \label{le:boundary}
\end{equation}
The sum is an even convex function of $x-X$ in $|x-X|<A$.
Thus on the circle it is at most
\begin{equation}
 \beta_\Gamma(\tau)=\min\left\{1,
 \operatorname{erfc}\frac{A-r}{2\sqrt\tau}
 +\operatorname{erfc}\frac{A+r}{2\sqrt\tau}\right\},
 \qquad \beta_\Gamma(0)=0.
 \label{le:beta}
\end{equation}

If $|f_y|\le\varepsilon_0$, $|g_\pm|\le M_b$ and
$|\mathcal R_N(t_*-v,\xi+\ii y)|\le\rho(v)$ over their full domains,
positivity gives the paid pointwise bound
\begin{equation}
 |E(u,x+\ii y)|\le\varepsilon_0S(u,x)+M_b(1-S(u,x))
                +\int_0^u\rho(v)S(u-v,x)\dd v.
 \label{le:mass-bound}
\end{equation}
In particular, when $M_b\ge\varepsilon_0$,
\begin{equation}
 \sup_\Gamma|E(u,z)|\le\varepsilon_0+
          (M_b-\varepsilon_0)\beta_\Gamma(u)+\int_0^u\rho(v)\dd v.
 \label{le:scalar-heat}
\end{equation}
For $M_b<\varepsilon_0$ keep \eqref{le:mass-bound}, or discard the
negative boundary correction. Inserting an upper bound for $1-S$ into
a negative multiple would be invalid. Spatially varying data may instead
be integrated in \eqref{le:bounded-duhamel}; the scalar suprema are only
one relaxation. The equally informed direct calculation receives this
propagated envelope too, and can intersect it with \eqref{le:source-value}.

\subsection{Recenter the logarithm before bounding the divisor}

The next elementary estimate concerns the full value-error disk, not
the quadrature aliases treated in the parent manuscript's common-log
Parseval estimate \cite[\S22]{VCDCMManuscript}. It uses the same principle
of retaining shared Fourier information.

\begin{lemma}[A centered logarithm disk]\label{le:log-disk}
For $0\le q<1$ and $|w|\le q$, define
\[
 \lambda=-\log(1-q),\qquad
 \kappa=\tfrac12\log(1-q^2),\qquad
 \eta=\operatorname{arctanh}q.
\]
The principal logarithm obeys
$|\log(1+w)-\kappa|\le\eta$. For $q>0$, $\eta<\lambda$.
\end{lemma}
\begin{proof}
The case $q=0$ is immediate. Put $v=(1+w)/\sqrt{1-q^2}$.
Its disk has center $\cosh\eta$ and radius $\sinh\eta$.
The Cayley map $\xi=(v-1)/(v+1)$ sends this disk onto
$|\xi|\le\tanh(\eta/2)$. The identity
$\log v=2\operatorname{arctanh}\xi$, on the branch containing $v=1$,
and the absolutely convergent odd power series give
$|\log v|\le2\operatorname{arctanh}|\xi|\le\eta$.
The disk lies in the right half-plane, so this is precisely the
principal logarithm in the statement. Finally,
$\lambda-\eta=-\tfrac12\log(1-q^2)>0$.
\end{proof}

\begin{theorem}[A complete coupled circle allowance]\label{le:gap-bound}
Let $H,P$ be holomorphic on a neighborhood of $|z-X|\le r$, have
exactly two zeros there counted with multiplicity, and be nonzero on
the circle. Suppose $|H/P-1|\le q<1$ on it. Define $c_P,D_P,c_H,D_H$
by the contour moments, and put $d=|c_P-X|/r$. Then
\begin{equation}
 |D_H-D_P|\le r^2 F(\eta,d),\qquad
 F(\eta,d)=\max_{0\le x\le\eta}
       \left\{\sqrt{\eta^2-x^2}+dx+\frac{x^2}{4}\right\}.
 \label{le:F}
\end{equation}
Here $F(0,d)=0$. Useful explicit upper bounds are
\begin{align}
 F(\eta,d)&\le\eta\sqrt{1+d^2}+\eta^2/4,
                                                        \label{le:cauchy}\\
 F(\eta,d)&\le\eta+\frac{\eta d^2}{2-\eta}\quad(0<\eta<2).
                                                        \label{le:parabolic}
\end{align}
For the centered subcase,
\begin{equation}
 F(\eta,0)=
 \begin{cases}\eta,&0\le\eta\le2,\\1+\eta^2/4,&\eta\ge2.
 \end{cases}
 \label{le:centered}
\end{equation}
The maximum $F$ optimizes the displayed two-coefficient relaxation;
the Cauchy and parabolic bounds are further upper relaxations. No
sharpness over physical heat fields is asserted.
\end{theorem}
\begin{proof}
Let $\ell=\log(H/P)$ on the circle collar and write
$\ell(X+re^{\ii\theta})-\kappa=\sum_{j\in\mathbb Z}a_j e^{\ii j\theta}$.
The strict relative clearance supplies the collar branch. Bessel's
inequality and Lemma~\ref{le:log-disk} give
$|a_{-1}|^2+|a_{-2}|^2\le\eta^2$. Constants have zero contour moments.
Substituting the two Fourier coefficients into
\eqref{prog:centered-error} gives exactly
\begin{equation}
 D_H-D_P=-r^2a_{-2}+(c_P-X)r a_{-1}-\frac{r^2}{4}a_{-1}^2.
 \label{le:two-modes}
\end{equation}
The square is a complex square, not a squared modulus. Taking absolute
values only after this identity, and putting $x=|a_{-1}|$, proves
\eqref{le:F}. To check attainment in the abstract coefficient ball,
write $(c_P-X)/r=de^{\ii\phi}$ and choose
$a_{-1}=-xe^{\ii\phi}$ and
$a_{-2}=\sqrt{\eta^2-x^2}\,e^{2\ii\phi}$; all three terms then
have the same phase. For $d=0$ any $\phi$ works.
Cauchy--Schwarz bounds the first two terms of its maximum
by $\eta\sqrt{1+d^2}$, proving \eqref{le:cauchy}.
Also $\sqrt{\eta^2-x^2}\le\eta-x^2/(2\eta)$.
Maximizing the resulting concave quadratic over the real line proves
\eqref{le:parabolic}. When $d=0$, set $y=x^2$ and maximize
$\sqrt{\eta^2-y}+y/4$ on $[0,\eta^2]$.
Its maximum is at $y=0$ for $\eta\le2$ and at $y=\eta^2-4$ otherwise,
which proves \eqref{le:centered}.
\end{proof}

The same source and geometry previously give the convenient bound
\begin{equation}
 B_{\rm old}=r^2\{\lambda\sqrt{1+d^2}+\lambda^2/4\}.
 \label{le:old}
\end{equation}
To see this directly, apply $|\ell|\le\lambda$ to the linear contour
moment, use the angular Cauchy--Schwarz bound
$\langle|re^{\ii\theta}+X-c_P|\rangle_\theta
\le r\sqrt{1+d^2}$, and retain the complete center-square payment.
For $q>0$, \eqref{le:cauchy} is strictly below \eqref{le:old}, because
$\eta<\lambda$. This is a strict improvement of that specified complete
allowance; it is not a comparison with an optimal direct estimator.
At $q=1/2,d=0$ the new allowance is $r^2\log(3)/2$, whereas
\eqref{le:old} is $r^2(\log2+(\log2)^2/4)$.
The ratio is less than $17/25$: more than $32\%$ of this full allowance
is saved. For an exact check, the positive $\operatorname{arctanh}$
series gives $\log2>56/81>69/100$ and
\[
 \tfrac12\log3
 =\operatorname{arctanh}(1/2)
 <\frac12+\frac1{24}+\frac1{160}
       +\frac1{7}\sum_{n=3}^{\infty}2^{-(2n+1)}
 =\frac12+\frac1{24}+\frac1{160}+\frac1{672}<\frac{11}{20}.
\]
Meanwhile $(17/25)(69/100+(69/100)^2/4)=0.550137>11/20$.
This finite constant example is not a canonical error sample.

\subsection{What the improved allowance does and does not protect}

For real conjugation-symmetric fields and a fixed counted circle, a
validated lower bound
\begin{equation}
 D_P(u)-r^2F(\eta(u),d(u))>0\quad\hbox{for every }0\le u\le U
 \label{le:margin}
\end{equation}
protects the enclosed real pair. Every error used to acquire $q$, the
count, $c_P$ and $D_P$ must be included. An endpoint instance alone
does not supply \eqref{le:margin}. Nor may the centered subcase be used
by silently moving the fixed circle to $c_P(u)$; the offset $d(u)$ pays
that distinction on the original contour.

If $H/P$ extends as a nonvanishing holomorphic function through the
interior, its logarithm is holomorphic there and both negative modes
vanish. The exact gap error is zero, even though a separate bound on the
joint second-derivative quotient need not vanish for a spatially varying
multiplier. In the full signed-time readout the center and heat-defect
terms supply the cancellation. Integrating that full identity gives
exactly the difference of the endpoint functional
\eqref{prog:centered-error}. Thus identical admissible source paths give
identical exact ranges for full transport and direct endpoint evaluation.
This is a scoped representation equality. Time-correlated data or a
useful heat acquisition may sharpen either calculation; no cost or
universal no-go claim follows.

For the exact control $H=P=(z-X)^2-2(t-T)$, $q=0$ and all new allowances
vanish, yet $D_P=2(t-T)$ reaches zero at its collision. The strict margin
\eqref{le:margin} fails. The remaining controls in Section~\ref{sec:program}
retain the generic heat/logarithm estimates with their own data. Changed
theta coefficients prevent importing \eqref{le:source-value} unchanged;
changed observation kernels fail the physical heat equation. Naming these
source failures does not establish a positive arithmetic margin. The
allowance improvement itself is generic and does not distinguish theta
from a colliding control.

\subsection{A finite arithmetic comparison inside known coverage}
\label{le:pilot}

The bounded pilot fixes
\begin{equation}
 X=140101/10,\quad A=2,\quad h=1,\quad r=9/10,\quad
 N=33,\quad 9/100\le t\le1/10.
 \label{le:pilot-cell}
\end{equation}
All fields and errors are multiplied by the same positive constant
$\exp(\pi X/8)$, independent of $z,t$. It changes neither the physical
heat equation nor the relative error or divisor. The cutoff inequalities
are verified on the entire rectangle. The final calculation uses
192-bit Arb/Acb arithmetic, 64 closed time cells and 512 closed angular
arcs. It proves $\Re(P_N(z)/(z-X)^2)>0$ on the count circle and the
annular boundary circles of radii $4/5$ and $1$, after the common scaling.
Thus the reference has count two and a single annular logarithm throughout
the time interval. The full source budget \eqref{le:source-value} is
included, and $q$ is the maximum of the local source/denominator ratios
on the \emph{same} arcs. Independent global suprema give a weaker bound
and are not used on one side of the comparison alone.

The reference moments are evaluated by the 512-node Laurent trapezoid
rule. If $M_-,M_+$ bound the annular logarithm on $r_-=4/5,r_+=1$,
the paid alias allowance for the weighted coefficient of order $k=1,2$ is
\begin{equation}
 M_+ r_+^k\frac{(r/r_+)^{512}}{1-(r/r_+)^{512}}
 +M_- r_-^k\frac{(r_-/r)^{512}}{1-(r_-/r)^{512}}.
 \label{le:aliases}
\end{equation}
This follows by bounding the positive and negative Laurent coefficient
ladders separately. Both complete allowances include these aliases,
interval time variation, reference-center uncertainty and arithmetic.
The largest coefficient alias is below $5.787\cdot10^{-23}$.
No exterior zero sum is truncated; the full field defines the contour
moments.

\begin{center}
\begin{tabular}{@{}lr@{}}
\toprule
Quantity on the entire pilot interval & Outward summary \\
\midrule
Complete relative-error allowance $q$ & $<0.166454$ \\
Reference discriminant lower bound & $>0.1418896$ \\
E10 transfer allowance \eqref{le:old} & $<0.1545636$ \\
Refined transfer allowance & $<0.1364749$ \\
Worst E10 complete lower bound (approximately) & $-0.0126739660$ \\
Refined complete all-time lower bound & $>0.0054147$ \\
\bottomrule
\end{tabular}
\end{center}

The comparison uses identical acquired $q,c_P,D_P$, source data and
numerical payments. The old sufficient lower bound is negative; this
says nothing adverse about the actual zeros. The new bound is positive
on every time cell and so protects this counted pair. Its transfer
allowance is smaller by more than $11\%$ in the worst-cell summary.
The separate centered $q=1/2$ example's $32\%$ is not the percentage of
this arithmetic pilot. The code, complete small record, method and
nonauthor replay are saved in the investigation's numerics, notes and
reviews folders. They constitute internal certificate checking with a
published source theorem as input, not independent specialist review.

This is a finite improvement of the specified direct estimator inside
an already verified canonical region. It gives no new height/time
coverage and no improved bound on the de Bruijn--Newman constant. Full
signed transport equals the same refined direct endpoint consumer on
identical information.

The scalar heat acquisition also has a definite scoped outcome. In the
same scaled units, its paid uniform bound is above $327656$, while the
raw direct contour allowance is below $196588$. The initial supremum is
already above $327394$; the lateral boundary probability at $U=1/100$
is below $7.358\cdot10^{-15}$. Thus this particular separate-supremum
heat estimate loses to the direct envelope. It does not show that the
actual residual is large, or exclude spatially weighted, signed or
correction-aware acquisitions. The two predeclared refinements and the
failed earlier enclosure are retained in the execution record.
